We implement constrained generation for leaf sections as a layered pipeline that (i) enforces a strict JSON output schema at decode time, (ii) embeds evidence-awareness into prompts and decoding constraints to discourage unsupported assertions, and (iii) applies automated post-generation validation and repair before a leaf section is accepted for assembly.
This design foregrounds auditability: every non-structural assertion produced for a leaf node includes a provenance field that either points to retrieved items or is explicitly marked as speculative for human review.
Decoding constraints and prompt-level clauses are provided to generators as a compact JSON schema specifying required fields (title, claim_list, paragraphs, citations, provenance) and type/length bounds.
Prompts include grounding clauses that require the model to (a) prefer wording traceable to provided retrieval snippets, (b) avoid introducing novel factual claims unless labeled as conjecture, and (c) populate citation slots with one or more source pointers drawn from the per-node retrieval set.
These prompt-level directives are combined with constrained decoding techniques, such as constrained-output wrappers or deterministic post-processing that accept only syntactically valid JSON.
Prior work surveys problems and practices for constrained neural generation and motivates careful constraint handling \cite{web_garbacea_2022_why_is_constrained_neural_language_generation_pa},
while recent analyses document limitations of constrained auto-regressive decoding and motivate conservative, bounded regeneration strategies in practice \cite{web_wu_2025_constrained_auto_regressive_decoding_constrains_}.
Automated validation runs immediately after generation and checks structural conformance (required keys, types, length bounds), citation integrity, and span-alignment (citation indices mapping to annotated token spans).
When a citation slot is empty but a supporting retrieval item matches a textual claim by conservative heuristics (exact substring match, anchor-phrase overlap), the pipeline attempts deterministic forced insertion of the matching retrieval identifier.
If no matching retrieval is found, the claim receives a standardized provenance tag (e.g., "no_evidence") so downstream tools and reviewers can identify and triage unsupported assertions.
Failed validations trigger a repair loop: deterministic fixes (JSON canonicalization, escape fixing, bracket balancing, and causal slot filling via exact-match heuristics) are attempted first.
If deterministic repairs fail, the system issues a constrained regeneration request returning the original retrieval snippets and validator diagnostics, instructing the model to prioritize exact quoting or close paraphrase tied to those snippets.
Regenerations are bounded by a small, configurable number of attempts; persistent failures escalate the leaf node to manual editing to preserve assembly correctness, a policy informed by observed decoder failure modes \cite{web_wu_2025_constrained_auto_regressive_decoding_constrains_}.
To minimize hallucination we enforce evidence-aware content rules: require explicit citation for factual assertions beyond high-level background, prefer short citation-linked claims over long unreferenced passages, and surface provenance metadata (retrieval ids, snippet offsets, confidence scores).
We also use dynamic retrieval patterns that can re-issue or refine retrievals conditioned on partially generated content, a strategy related to context-aware retrieval during generation \cite{web_qi_2025_ar_rag_autoregressive_retrieval_augmentation_for}; any speculative assertion must include an explicit rationale and reviewer guidance.
Sections with speculative labels, failed repairs, or low citation coverage are routed to lightweight human review, where edits are applied to the JSON, re-validated, and recorded in an audit trail (timestamps, editor id, before/after states) to preserve traceability.